This chapter synthesises the per-requirement evaluation results from Chapter~\ref{chap:evaluation} into a cross-cutting analysis, positions the measured outcomes against the related work surveyed in Chapter~\ref{chap:analysis}, and discusses the limitations that constrain the current findings. The goal is not to repeat the per-requirement numbers but to interpret them.

\section{Result Analysis}

The strongest signal across the evaluated configurations is the gap between an obvious-sink category and a semantic category. Code Guardian's headline configuration (\texttt{qwen3:8b+RAG}) achieves 100\% canonical recall on command injection, hardcoded credentials, header injection, information exposure, LDAP injection, path traversal, prototype pollution, SSRF, weak cryptography, and XSS, and 90\% on SQL injection. The same configuration achieves 0\% recall on improper authentication, input validation, and weak encryption. The qualitative analysis in Section~\ref{sec:eval-error-analysis} shows that the improper-auth and weak-encryption gaps are taxonomy artefacts (the model emits sibling canonical categories for the same code), while input-validation is a genuine reasoning gap on real-world library code. This split is the operational story of the thesis: contextual reasoning about lexically-explicit sinks works very well; reasoning about the absence of validation does not.

False-positive control improved decisively. The previously-published 20\% sample-level FPR on \texttt{qwen3:8b+RAG} drops to 10\% under the rerun's symmetric three-runs-per-secure-sample configuration, and the inter-run confidence gate at $\geq 0.67$ keeps that 10\% while pushing precision from 73.53\% to 77.78\%. Most other LLM configurations still flag every secure case at least once, even with the gate, which means the control is not a generic property of multi-pass consensus --- it is specific to the strongest model in the set. This is consistent with the broader observation in the literature that small models trade calibration for speed.

Latency results validate the design hypothesis that prefill-token cost dominates inference on local hardware. The tightened system prompt cut median latency on \texttt{qwen3:8b+RAG} from 2\,721\,ms to 1\,573\,ms ($-42\%$) without changing the model, hardware, or generation parameters. Stage-split telemetry confirms inference is more than 99\% of the total response time; there is no realistic gain to be made on the parsing side. The remaining latency budget is therefore best spent on either smaller models (with a quality cost) or on architectural changes such as streaming partial JSON output (deferred to future work).

The repair-quality result is the most surprising negative finding. Across all 297 generated repair attempts in the \texttt{qwen3:8b+RAG} configuration, zero passed the syntactic auto-applicability check. The model fills the \texttt{suggestedFix} string with prose openings (``Use parameterized queries\ldots'', ``Avoid string concatenation\ldots'') even though the JSON response schema constrains the field type to a string. The schema constrains the field type, not the content language. This is a contract-design problem rather than a model-quality problem.

\section{Comparison with Related Work}

Three benchmarks anchor the comparison. IRIS~\cite{li2025iris} reports a false discovery rate of 84.82\% (precision approximately 15\%) for GPT-4 on whole-repository Java analysis, with 55 of 120 vulnerabilities detected. CWEval~\cite{peng2025cweval} establishes outcome-driven evaluation as a methodological benchmark and reports that LLMs achieve more than 70\% F1 with about 80\% precision when given context-rich prompts. SecRepair~\cite{islam2024secrepair} reports a 12\% improvement on automated similarity-based repair-quality metrics but explicitly admits in its own analysis that those metrics are ``fundamentally inadequate'' for security correctness assessment.

Code Guardian's measured 73.53\% precision on a local 8B-parameter model lands in the same precision class as the cloud-LLM systems cited above. This is not a head-to-head dominance claim --- the datasets and analysis scopes differ --- but it is a defensible counterpoint to the assumption that local models cannot reach cloud-class precision on focused vulnerability detection. The path that gets there is the explicit pipeline of canonical type matching, multi-pass consensus, inter-run confidence gating, and a tightened structured-output contract, rather than raw model scale.

The honest negative result on RAG also aligns with the literature. RESCUE~\cite{shi2025rescue} reports that conventional retrieval augmentation (SecCoder, the closest baseline in their study) does not significantly improve secure code generation, and proposes a more sophisticated multi-faceted retrieval to close the gap. The thesis's RAG ablation results show \texttt{qwen3:8b} benefitting (+4.42 F1 with the canonical matcher) while \texttt{codellama} collapses ($-29.89$ F1 with RAG enabled). The collapse is driven by the additional retrieval context interfering with \texttt{codellama}'s output structure: latency nearly doubles, recall drops by twenty points, and parse failures increase. The conclusion that RAG must be calibrated per model rather than enabled by default is therefore strongly supported, both by the rerun data and by the related literature.

The 0\% auto-applicable repair rate is consistent with SecRepair's own characterisation of automated repair-quality metrics. Where SecRepair concludes that those metrics are inadequate without quantifying the inadequacy, this thesis quantifies it: zero of 297 generated repairs satisfy a strict syntactic-applicability bar, even with a JSON schema in place. The auto-applicable rate is therefore positioned as a deterministic, fleet-wide alternative to similarity-based repair metrics, addressing one specific axis of the inadequacy SecRepair identifies.

\section{Deviations from the Approved Task}

A subset of mechanisms named in the registered task description (Page~iv) was scoped down or replaced during execution; the deviations are recorded here to keep the delivered work and the approved task aligned.

\paragraph{Datasets: benchmark suites.}
The approved task specified evaluation on \emph{Juliet} and \emph{OWASP Benchmark} (40--50 CWE-mapped cases). Both suites are real but neither targets JavaScript: Juliet is published by NIST in C/C++ and Java only, and the OWASP Benchmark Project is a Java application. There is no first-party JavaScript port of either, and running this thesis's JS/TS detector against C/C++ or Java would not be meaningful. The delivered evaluation therefore substitutes two CWE-mapped JavaScript corpora, both documented in Section~\ref{sec:eval-benchmark-substitution}: the curated 101-case primary corpus (71 vulnerable + 30 secure) and an additional 15-case external corpus drawn from OWASP NodeGoat, OWASP Juice Shop, and three named CVEs. Per-case provenance is recorded in \texttt{evaluation/datasets/SOURCES.md}. The substitution preserves the property the named benchmarks provide --- CWE-mapped, vulnerability-tagged samples drawn from authoritative sources --- in the language the detector targets. The substitution penalty is measured in Section~\ref{sec:eval-r1-external-corpus}: under deliberately weak conditions (LLM-only, single run, no consensus, no confidence gate) the external corpus produced F1 40.00\% versus the curated headline's 71.94\%, with most misses in framework-level weakness classes (mass assignment, JWT \texttt{none}-algorithm acceptance, signature verification, Lodash template injection) that the function-level scope does not see.

\paragraph{Metrics.}
Precision, recall, F1-score, and latency are delivered as specified. Resource usage is reported per-inference in every run JSON via \texttt{process.cpuUsage()} and sampled \texttt{process.memoryUsage()} (Section~\ref{sec:eval-r4-resources}). Privacy validation is delivered as both an architectural checklist (Section~\ref{sec:eval-r5}) and an automated egress-and-leak harness (\texttt{evaluation/privacy/}, Section~\ref{sec:eval-r5-egress}). Reproducibility is delivered as the per-run \texttt{byteIdentityRate} and \texttt{setIdentityRate} measured by \texttt{evaluation/reproducibility.js} on a fixed-seed re-run subset (Section~\ref{sec:eval-reproducibility-metric}). The harness additionally reports JSON-parse success rate, secure-sample false-positive rate at sample level, and the auto-applicable repair rate; these were added during execution to give a fuller picture of structured-output reliability and repair quality.

\paragraph{Real-world project analysis.}
The approved task names ``one actively maintained real-world JavaScript/TypeScript project with documented vulnerabilities and verified patches''. The delivered evaluation runs the detector against OWASP NodeGoat at the \texttt{master} commit via \texttt{evaluation/run-realworld.js} (Section~\ref{sec:eval-realworld}). The result is recorded in Section~\ref{sec:eval-r1-realworld}: 26 source files scanned, 85 function chunks analysed, 14 plausible findings produced, and 0/10 recall against the curated 10-entry OWASP-Top-10 documented-vulnerability list. The 0\,\% recall is honestly disclosed and decomposed into four mechanisms (out-of-scope file types, cross-file flows the function-level scope cannot see, canonical-taxonomy under-coverage, and two genuine LLM-side false negatives in \texttt{data/allocations-dao.js} and \texttt{data/profile-dao.js}). The run validates the thesis's existing function-scope limitation rather than refuting the headline numbers: the detector finds plausible real weaknesses on real code, but recall against a human-curated vulnerability list is bounded by what a function-level scope can see. Cross-file context modelling and canonical-taxonomy expansion are named as future work.

These deviations were discussed with the supervisor during the implementation phase.

\section{Limitations}

\paragraph{Scope and context depth.}
The system handles common vulnerability patterns but remains limited for deep cross-file reasoning without full static data-flow analysis. A regex-based import / sink resolver (Section~\ref{sec:impl-tech-stack}) provides lightweight cross-snippet hints, but its empirical effect on the four semantic categories that previously scored zero recall was marginal --- the canonical taxonomy recovers most of the apparent gap by re-binning labels rather than by detecting new instances.

\paragraph{Evaluation representativeness.}
The curated dataset is intentionally small and human-auditable. The 101 cases align with the scale of contemporary curated benchmarks (CWEval at 119 tasks, IRIS at 120 vulnerabilities) but do not generalise to production repository scale; results are indicative, not definitive.

\paragraph{Repair correctness bar.}
Repairs are model-generated and may change behaviour. Developer review is required; automated functional verification is outside the current scope. The pipeline includes a syntactic auto-applicability check (Section~\ref{sec:eval-r3-auto-applicable}) that scales to the full evaluation set, but this is a strictly weaker signal than functional verification: a parse-clean fix can still be semantically wrong, and a parse-failing fix may still be informative as decision support.

\paragraph{Hardware and run-to-run variance.}
All measurements are from a single Apple M4 Max profile (16 cores, 64\,GB RAM) with deterministic decoding (\texttt{temperature=0}, fixed seed). Local inference on GPU hardware can introduce minor non-determinism from floating-point scheduling. The 3-run design partially mitigates this for vulnerable cases, but performance on lower-end hardware may differ. Reported values are descriptive measurements under a documented environment, not hardware-independent constants.

\paragraph{Single-rater R3 manual review.}
The 25-sample manual review of repair quality (Section~\ref{sec:eval-r3}) was performed by a single reviewer; no inter-rater reliability was measured. The deterministic auto-applicable rate complements this on the full $n=297$ sample but is a stricter and narrower bar.

\paragraph{No user study for R4.}
Usability is measured by latency only. Developer satisfaction, trust, and adoption were not evaluated. The R4 thresholds reflect HCI response-time guidelines rather than empirical usability data on this specific tool.

\section{Summary}

The cross-requirement analysis shows a coherent picture: local 8B-parameter models reach cloud-LLM-class precision on focused vulnerability detection when paired with explicit pipeline mechanisms (canonical type matching, multi-pass consensus, inter-run confidence gating, tight structured-output contract); they remain weaker on semantic categories that require reasoning about absent code; and the schema-only repair contract caps the achievable auto-applicable repair rate. The next chapter draws these threads together into a final summary and outlook.
